% Section 2 - The standard library
% Roberto Masocco <roberto.masocco@uniroma2.it>
% April 22, 2026

% ### The standard library ###
\section{The standard library}
\graphicspath{{figs/section2/}}

% --- The standard library ---
\begin{frame}{The standard library}
	\justifying
	The \textbg{C++ Standard Library} is a large collection of ready-made components, all living in the \texttt{std} namespace.\\
	It is built almost entirely on templates, and provides
	\begin{itemize}
		\item \textbg{containers}: data structures for storing and organizing objects;
		\item \textbg{algorithms}: sorting, searching, transforming, accumulating, etc.;
		\item \textbg{utilities}: string handling, time, random numbers, type traits, etc.;
		\item \textbg{I/O streams}: file and console I/O;
		\item \textbg{concurrency}: threads, mutexes, atomics, condition variables.
	\end{itemize}
	\bigskip
	For this course, a working knowledge of a small but important subset is sufficient.\\
	The rest you will pick up as needed — the documentation is excellent.
\end{frame}
\begin{frame}[fragile]{The standard library}{\texttt{std::string}}
	\justifying
	It manages its own memory, supports concatenation with \texttt{+}, comparison with \texttt{==}, and all the operations you would expect; \textbg{conversion} to \texttt{const char*} is available via \texttt{s.c\_str()}.
	\begin{columns}
		\column{.9\textwidth}
		\begin{lstlisting}[style=codebeamer, language=C++, caption=Basic \texttt{std::string} usage.]
#include <string>

std::string s1 = "hello";
std::string s2 = "world";
std::string s3 = s1 + ", " + s2 + "!"; // "hello, world!"

std::cout << s3 << std::endl;
std::cout << s3.size() << std::endl;        // 13
std::cout << s3.substr(0, 5) << std::endl;  // "hello"

if (s1 == "hello") { /* ... */ }       // value comparison
    \end{lstlisting}
	\end{columns}
\end{frame}
\begin{frame}{The standard library}{\texttt{std::vector}}
	\justifying
	\texttt{std::vector<T>} is a \textbg{dynamically-sized array} — the workhorse container of C++.\\
  \bigskip

	It stores elements \textbg{contiguously} in heap memory, \textbg{grows automatically} when needed, and provides O(1) indexed access.\\
  \bigskip
	
	Prefer \texttt{push\_back} and indexed access for most uses.\\
	When the final size is known in advance, \texttt{v.reserve(n)} avoids repeated \textbg{reallocations}.
\end{frame}
\begin{frame}[fragile]{The standard library}{\texttt{std::vector}}
  \begin{columns}
		\column{.9\textwidth}
		\begin{lstlisting}[style=codebeamer, language=C++, caption=Basic \texttt{std::vector} operations.]
#include <vector>

std::vector<int> v;
v.push_back(10);          // append
v.push_back(20);
v.push_back(30);

int first = v[0];         // 10, no bounds check
int safe  = v.at(1);      // 20, throws if out of range
std::size_t n = v.size(); // 3

v.pop_back();             // removes last element
    \end{lstlisting}
	\end{columns}
\end{frame}
\begin{frame}{The standard library}{Iterating containers}
	\justifying
	C++ containers expose \textbg{iterators} — objects that behave like pointers and allow traversal.\\
	Writing iterator types by hand is verbose; two modern features eliminate the boilerplate.\\
	\bigskip
	\texttt{auto} asks the compiler to \textbg{deduce the type} of a variable from its initializer — useful whenever the type is obvious from context or too verbose to write.\\
	\bigskip
	The \textbg{range-based} \texttt{for} loop iterates over any container without index arithmetic.
\end{frame}
\begin{frame}[fragile]{The standard library}{Iterating containers}
  \begin{columns}
		\column{.9\textwidth}
		\begin{lstlisting}[style=codebeamer, language=C++, caption=Iteration with \texttt{auto} and range-based \texttt{for}.]
std::vector<int> v = {1, 2, 3, 4, 5};

for (auto x : v)           // x is a copy of each element
  std::cout << x << " ";

for (auto & x : v)         // x is a reference, no copy
  x *= 2;                  // modifies the vector in place

for (const auto & x : v)   // read-only reference
  std::cout << x << " ";
    \end{lstlisting}
	\end{columns}
\end{frame}
\begin{frame}{The standard library}{Other essential containers}
	\justifying
	\texttt{std::array<T, N>} is a \textbg{fixed-size array} whose size is known at compile time.\\
	Unlike a raw C array it knows its own size, can be copied, and works with standard algorithms.\\
	\bigskip
	\texttt{std::map<K, V>} is an \textbg{ordered associative container}: a sorted tree of key-value pairs with O(log n) lookup.
\end{frame}
\begin{frame}[fragile]{The standard library}{Other essential containers}
  \begin{columns}
		\column{.9\textwidth}
		\begin{lstlisting}[style=codebeamer, language=C++, caption=\texttt{std::array} and \texttt{std::map} basics.]
#include <array>
#include <map>

std::array<double, 3> point = {1.0, 2.0, 3.0};
std::cout << point[2] << std::endl;    // 3.0

std::map<std::string, int> scores;
scores["alice"] = 95;
scores["bob"]   = 87;

for (const auto & [key, val] : scores) // structured binding
  std::cout << key << ": " << val << std::endl;
    \end{lstlisting}
	\end{columns}
\end{frame}
